\begin{helper}
Let \(g:\mathbb R^3\to\mathbb R\) be continuous.  Then the all-real iterated
nonnegative integral of the simplex indicator equals the nested
nonnegative integral over the ordered simplex:
\[
\int_{\mathbb R}^{+}\int_{\mathbb R}^{+}\int_{\mathbb R}^{+}
\mathbf 1_{\{0\le x\le y\le z\le1\}}
\operatorname{ofReal}(g(x,y,z))\,dz\,dy\,dx
=
\int_{(0,1]}^{+}\int_{(0,z]}^{+}\int_{(0,y]}^{+}
\operatorname{ofReal}(g(x,y,z))\,dx\,dy\,dz .
\]
\end{helper}
\begin{proof}
Let
\[
S=\{(x,y,z):0\le x\le y\le z\le1\}.
\]
The set \(S\) is Borel measurable, since it is the finite intersection of
closed sets defined by the continuous inequalities \(x\ge0\), \(y-x\ge0\),
\(z-y\ge0\), and \(1-z\ge0\).  Continuity of \(g\) implies measurability of
\((x,y,z)\mapsto\operatorname{ofReal}(g(x,y,z))\), so the indicator integrand
is a measurable \([0,\infty]\)-valued function.

Tonelli's theorem applies to this nonnegative measurable function on
\(\mathbb R^3\).  It permits changing the order of integration from the
displayed \(x\)-then-\(y\)-then-\(z\) all-real iterated integral to the
\(z\)-then-\(y\)-then-\(x\) all-real iterated integral without changing the
value.

For fixed \(z\), the indicator is zero unless \(0\le z\le1\).  Replacing the
outer all-real integral by the restricted integral over \((0,1]\) therefore
does not change the value, because the discarded boundary point \(z=0\) has
Lebesgue measure zero and the integrand vanishes off \([0,1]\).  For such a
fixed \(z\), the indicator is zero unless \(0\le y\le z\), so the middle
all-real integral restricts to \((0,z]\), again with only a null endpoint
discarded.  For fixed \(y,z\) in those ranges, the indicator is zero unless
\(0\le x\le y\), so the inner all-real integral restricts to \((0,y]\).  On
the remaining region \(0<x\le y\le z\le1\), the indicator is identically one,
and the integrand is exactly \(\operatorname{ofReal}(g(x,y,z))\).  These three
restriction steps give the displayed equality.
\end{proof}
